\section{The raw-graph binding certificate}

The abstract theorem does not assign receipts to an unlabeled graph. Applying
the algebra requires a separate binding certificate whose inputs are derived
without consulting a desired holonomy or component answer.

\subsection*{Structural inputs}

A complete certificate must identify and hash-lock:

\begin{enumerate}
\item a connected base graph $X$ with its dart-reversal convention;
\item a candidate lifted graph $Y$ and projection $p:Y\to X$;
\item a finite group $R$ with an explicit multiplication table;
\item a free and transitive $R$-action on every fiber of $p$;
\item an edge-receipt assignment $\alpha:D(X)\to R$;
\item the inverse law $\alpha(\bar a)=\alpha(a)^{-1}$.
\end{enumerate}

Equal fiber size alone does not determine the receipt group. A four-element
fiber, for example, does not decide between $C_4$, $V_4$, or a set with no
admitted group structure. The regular action must be exhibited.

\subsection*{Reconstruction gate}

The certificate must reconstruct the predicted lift from $(X,R,\alpha)$ and
provide a projection-preserving, $R$-equivariant graph isomorphism
\[
Y\cong\widetilde X_{\alpha}.
\]
This prevents a receipt table from being accepted merely because its final
holonomy looks useful.

\subsection*{Invariant output}

After reconstruction, the packet should export:

\begin{enumerate}
\item a spanning tree and fundamental circuit basis;
\item each fundamental circuit receipt;
\item the generated subgroup $H_{\alpha}$;
\item predicted and observed component counts;
\item root-change and spanning-tree-change checks;
\item gauge covariance checks;
\item the exact claim boundary for orientation and physical interpretation.
\end{enumerate}

If an absolute direction is required, it must be supplied by independently
oriented source structure or a directed arriving action. The receipt algebra
classifies what has been bound; it does not manufacture the missing arrow.
